%================================================================
% APPENDIX: TYPE-A BAXTER--REES COMPACTIFICATION AND BOUNDARY
% TENSOR GEOMETRY
% Status: FRONTIER appendix. This appendix packages the thread's new
% mathematics in a repository-ready form. It records the weightwise
% completion formalism used in the MC4 proof and sharpens the DK
% frontier recorded in Chapter~\ref{chap:concordance} and in
% Appendices~\ref{app:nilpotent-completion} and~\ref{app:coderived}.
% Note: MC4 is now proved (thm:winfty-all-stages-rigidity-closure).
%================================================================

\chapter{Type-$A$ Baxter--Rees compactification, weightwise MC4, and boundary tensor geometry}
\label{app:typeA-baxter-rees-theta}

\begin{quote}
This appendix reorganizes the Yangian frontier from the Chriss--Ginzburg
standpoint of the main text (every algebraic structure is an MC element
in a convolution algebra) and records a theorem package tailored to the
type-$A$ derived Drinfeld--Kohno problem. The weight filtration on the
RTT convolution algebra separates three scales. At the compact scale, the
polynomial core generated by the vector evaluation line is rigid: one
binary seed on $V\otimes V$ propagates to all Schur functors, and in the
strict regime it forces all higher-point corrections. At the completion
scale, the principal RTT mode tower for $Y_\hbar(\mathfrak{sl}_N)$
satisfies \emph{weightwise} stabilization (not unweighted stabilization),
and this is the form in which MC4 is proved. At the noncompact scale,
asymptotic and negative prefundamental modules are constructed as derived
telescope fibers of Kirillov--Reshetikhin chains, assembled into a single
affine Baxter--Rees family over $\mathbb A^1_q$ whose generic fiber is the
ordinary asymptotic module and whose special fiber is the negative
prefundamental module.

Combining this Baxter--Rees family with the polynomial $R$-matrix theory
of shifted Yangians and with Zhang's Theta-associators, we extend objects
together with their braiding and associator data across the boundary,
yielding a canonical formal braided boundary germ with a boundary
Kodaira--Spencer class. The extension problem beyond the compact
polynomial core reduces to a finite set of explicit verification tasks.
\end{quote}

\begin{remark}[Post-closure note]
MC4 is now proved unconditionally via $W_N$ rigidity
(Theorem~\ref{thm:winfty-all-stages-rigidity-closure}). The weightwise
completion formalism developed in this appendix is part of the proof.
\end{remark}

\section{Organizing principle and constitutional role}
\label{sec:typeA-organizing-principle}

The concordance chapter isolates three live frontier packages for the
Yangian story: the factorization-level lift of the derived
Drinfeld--Kohno square, the infinite-generator completion boundary, and
the BV/BRST/bar comparison; see Chapter~\ref{chap:concordance}. The
completion appendices already explain that any genuine extension beyond
the compact finite-type regime must pass through completed and, in the
curved case, coderived categories; see
Appendices~\ref{app:nilpotent-completion} and~\ref{app:coderived}.
The present appendix sharpens that frontier in type~$A$.

The organizing principle is the Chriss--Ginzburg framework of the main
text: every algebraic structure is a Maurer--Cartan element in a
convolution algebra. The Rees-parameter weight grading on the RTT
presentation introduces a natural filtration on the convolution dg~Lie
algebra $\Convstr(C^!_{\mathrm{ch}},\, P^{\mathrm{ch}})$, and the key
observation is that this filtration makes completion automatic: the
inverse system stabilizes weightwise, so the completed bar complex is
the product of its weight-stable pieces rather than a nontrivial
pro-object. The three arrows
\begin{equation}
\text{compact seed} \quad \longrightarrow \quad
\text{weightwise completion} \quad \longrightarrow \quad
\text{formal tensor boundary germ}
\end{equation}
are therefore arrows \emph{within} the convolution algebra: the seed is
the binary MC element $r(z)$ on the polynomial core, the completion is
the weight-adic pro-MC lift, and the boundary germ is the formal
neighborhood of the MC element in the completed moduli space.

\begin{remark}[Status and scope]
\label{rem:typeA-status-scope}
This appendix contains three kinds of statements.
\begin{enumerate}
 \item Unconditional internal theorems whose proofs are given here,
 such as the seed propagation theorem, the weightwise MC4 theorem for
 the principal RTT tower, the telescope constructions, the
 Baxter--Rees family, and the continuation principles for polynomial
 tensor data.
 \item Conditional theorems whose hypotheses isolate the precise
 external input required from the literature, especially from the
 polynomial $R$-matrix theory of shifted Yangians and from Theta
 associators for completed triple tensor products
 \cite{HZ24,ZhangTheta24}.
 \item Reduction theorems showing that once the compact factorization
 equivalence is known on the Kirillov--Reshetikhin packet, the
 extension to the completed/coderived Baxter envelope becomes formal.
\end{enumerate}
The purpose is not to blur the distinction between theorematic core and
frontier input, but to compress the frontier into a smaller and sharper
set of verification problems.
\end{remark}

\section{Seed rigidity on the compact polynomial core}
\label{sec:typeA-seed-rigidity}

We begin with the smallest place in which the ordered Yangian structure
is already visible: the idempotent-complete tensor subcategory generated
by the vector representation and its spectral shifts.

\subsection{Additive spectral kernels}

Let $V=\mathbb C^N$ denote the standard $\mathfrak{gl}_N$-module. Write
$\mathcal C_{\mathrm{poly}}$ for the idempotent-complete tensor category
generated by $V$ under tensor products, direct summands, and spectral
shifts. An object of $\mathcal C_{\mathrm{poly}}$ is thus a finite direct
sum of Schur functors $S_\lambda V$ with spectral shifts suppressed from
the notation.

\begin{definition}[Additive spectral kernel]
\label{def:additive-spectral-kernel}
An \emph{additive spectral kernel} on $\mathcal C_{\mathrm{poly}}$ is a
family of meromorphic operators
\begin{equation}
r_{X,Y}(u)\in \operatorname{End}(X\otimes Y)((u^{-1})),
\qquad X,Y\in \mathcal C_{\mathrm{poly}},
\end{equation}
such that
\begin{align}
r_{\mathbf 1,X}(u)&=0=r_{X,\mathbf 1}(u),
\label{eq:additive-kernel-unit}\\
r_{X\otimes X',Y}(u)&=r_{X,Y}(u)^{13}+r_{X',Y}(u)^{23},
\label{eq:additive-kernel-left}\\
r_{X,Y\otimes Y'}(u)&=r_{X,Y}(u)^{12}+r_{X,Y'}(u)^{13}.
\label{eq:additive-kernel-right}
\end{align}
The superscripts indicate the usual legs in iterated tensor products.
\end{definition}

This is the strict binary shadow of the dg-shifted Yangian structure
appearing in the line-operator theory of
Dimofte--Niu--Py~\cite{DNP25}. In the quasi-linear regime of
\cite{DNP25}, the higher $A_\infty$-operations vanish and the binary
kernel becomes the whole story.

\begin{theorem}[Propagation from the vector seed]
\ClaimStatusProvedHere
\label{thm:vector-seed-propagation}
Let $r$ be an additive spectral kernel on
$\mathcal C_{\mathrm{poly}}$, and set
\begin{equation}
R(u):=r_{V,V}(u)\in \operatorname{End}(V^{\otimes 2})((u^{-1})).
\end{equation}
Then for all $m,n\ge 0$,
\begin{equation}
\label{eq:vector-seed-propagation}
r_{V^{\otimes m},V^{\otimes n}}(u)
=
\sum_{i=1}^m\sum_{j=1}^n R_{i,m+j}(u).
\end{equation}
In particular the entire additive kernel is uniquely determined by the
single binary seed $R(u)$.
\end{theorem}

\begin{proof}
We argue by induction on $m+n$. The cases $m=0$ or $n=0$ are exactly
\eqref{eq:additive-kernel-unit}. Assume the formula known for
$(m,n)$. Then
\begin{equation}
V^{\otimes(m+1)}=V^{\otimes m}\otimes V,
\end{equation}
so \eqref{eq:additive-kernel-left} gives
\begin{equation}
r_{V^{\otimes(m+1)},V^{\otimes n}}(u)
=
r_{V^{\otimes m},V^{\otimes n}}(u)^{13}
+
r_{V,V^{\otimes n}}(u)^{23}.
\end{equation}
By the induction hypothesis,
\begin{equation}
r_{V^{\otimes m},V^{\otimes n}}(u)^{13}
=
\sum_{i=1}^{m}\sum_{j=1}^{n}R_{i,m+1+j}(u).
\end{equation}
A second induction on $n$ using \eqref{eq:additive-kernel-right} shows
that
\begin{equation}
r_{V,V^{\otimes n}}(u)=\sum_{j=1}^n R_{1,1+j}(u),
\end{equation}
hence
\begin{equation}
r_{V,V^{\otimes n}}(u)^{23}
=
\sum_{j=1}^n R_{m+1,m+1+j}(u).
\end{equation}
Adding the two expressions yields \eqref{eq:vector-seed-propagation} for
$(m+1,n)$. The induction in the second variable is identical. Uniqueness
is immediate because every object of $\mathcal C_{\mathrm{poly}}$ is built
from $V$ by tensor products, direct summands, and finite direct sums.
\end{proof}

\subsection{The Schur-envelope formula}

\begin{lemma}[Equivariant normal form on $V\otimes V$]
\ClaimStatusProvedHere
\label{lem:equivariant-normal-form}
If $R(u)\in \operatorname{End}(V\otimes V)((u^{-1}))$ commutes with the
diagonal $\mathfrak{gl}_N$-action, then there exist unique scalar Laurent
series $f(u),g(u)\in \mathbb C((u^{-1}))$ such that
\begin{equation}
R(u)=f(u)P+g(u)\operatorname{id}_{V\otimes V},
\end{equation}
where $P(x\otimes y)=y\otimes x$ is the flip operator.
\end{lemma}

\begin{proof}
The $\mathfrak{gl}_N$-module $V\otimes V$ decomposes as
\begin{equation}
V\otimes V\cong S^2V\oplus \Lambda^2V,
\end{equation}
and both summands appear with multiplicity one. By Schur's lemma, any
$\mathfrak{gl}_N$-equivariant endomorphism acts by a scalar on each
summand. Writing the corresponding projectors as
\begin{equation}
\pi_s=\frac{\operatorname{id}+P}{2},
\qquad
\pi_a=\frac{\operatorname{id}-P}{2},
\end{equation}
we obtain
\begin{equation}
R(u)=a(u)\pi_s+b(u)\pi_a
=\frac{a(u)-b(u)}{2}P+\frac{a(u)+b(u)}{2}\operatorname{id},
\end{equation}
which is the desired form.
\end{proof}

For polynomial $\mathfrak{gl}_N$-modules $X,Y$ define the cross-Casimir
operator
\begin{equation}
\Omega_{X,Y}:=\sum_{a,b=1}^N \rho_X(E_{ab})\otimes \rho_Y(E_{ba}).
\end{equation}

\begin{theorem}[Schur-envelope formula]
\ClaimStatusProvedHere
\label{thm:schur-envelope-formula}
Let $r$ be an additive $\mathfrak{gl}_N$-equivariant spectral kernel with
vector seed
\begin{equation}
R(u)=f(u)P+g(u)\operatorname{id}_{V\otimes V}.
\end{equation}
Then for every pair of Schur functors $X=S_\lambda V$ and
$Y=S_\mu V$ one has
\begin{equation}
\label{eq:schur-envelope-formula}
r_{X,Y}(u)=f(u)\,\Omega_{X,Y}+g(u)\,|\lambda|\,|\mu|\,\operatorname{id}_{X\otimes Y}.
\end{equation}
\end{theorem}

\begin{proof}
By Theorem~\ref{thm:vector-seed-propagation}, for tensor powers we have
\begin{equation}
r_{V^{\otimes m},V^{\otimes n}}(u)
=
f(u)\sum_{i=1}^{m}\sum_{j=1}^{n}P_{i,m+j}
+g(u)mn\,\operatorname{id}.
\end{equation}
On the other hand,
\begin{equation}
\Omega_{V^{\otimes m},V^{\otimes n}}
=
\sum_{a,b}\left(\sum_{i=1}^{m}E_{ab}^{(i)}\right)
\left(\sum_{j=1}^{n}E_{ba}^{(m+j)}\right)
=
\sum_{i=1}^{m}\sum_{j=1}^{n} \sum_{a,b}E_{ab}^{(i)}E_{ba}^{(m+j)}.
\end{equation}
Since on $V\otimes V$ one has
\begin{equation}
\sum_{a,b}E_{ab}\otimes E_{ba}=P,
\end{equation}
it follows that
\begin{equation}
\Omega_{V^{\otimes m},V^{\otimes n}}=\sum_{i,j}P_{i,m+j}.
\end{equation}
Hence
\begin{equation}
r_{V^{\otimes m},V^{\otimes n}}(u)
=
f(u)\Omega_{V^{\otimes m},V^{\otimes n}}+g(u)mn\operatorname{id}.
\end{equation}
Now $S_\lambda V$ and $S_\mu V$ are direct summands of
$V^{\otimes |\lambda|}$ and $V^{\otimes |\mu|}$ cut out by Young
idempotents, and $\Omega$ commutes with the symmetric-group actions.
Restricting the tensor-power identity to the corresponding direct summands
gives \eqref{eq:schur-envelope-formula}.
\end{proof}

\begin{corollary}[One ordered coefficient forces the reduced core]
\ClaimStatusProvedHere
\label{cor:one-coefficient-forces-reduced-core}
Let $r^{\mathrm{red}}$ denote the image of $r$ modulo scalar multiples of
the identity on each $X\otimes Y$. Suppose that
\begin{equation}
\label{eq:ordered-coefficient-condition}
r^{\mathrm{red}}_{V,V}(u)(e_1\otimes e_2)
=
-\hbar u^{-1}e_2\otimes e_1.
\end{equation}
Then for every polynomial pair $X,Y\in \mathcal C_{\mathrm{poly}}$,
\begin{equation}
\label{eq:reduced-core-forced}
r^{\mathrm{red}}_{X,Y}(u)
=
-\hbar u^{-1}\Omega^{\mathfrak{sl}_N}_{X,Y}.
\end{equation}
\end{corollary}

\begin{proof}
By Lemma~\ref{lem:equivariant-normal-form}, the reduced seed on
$V\otimes V$ is determined by the coefficient of the flip channel. The
condition \eqref{eq:ordered-coefficient-condition} therefore forces
$f(u)=-\hbar u^{-1}$. Passing to the reduced quotient removes the scalar
term $g(u)\operatorname{id}$ in
Theorem~\ref{thm:schur-envelope-formula}. The result is exactly
\eqref{eq:reduced-core-forced}.
\end{proof}

The next statement upgrades binary rigidity to all ordered
configurations in the strict regime.

\begin{definition}[Strict binary regime]
\label{def:strict-binary-regime}
A dg-shifted Yangian is said to be in the \emph{strict binary regime} if
all higher $A_\infty$-operations $m_k$ vanish for $k\ge 3$, so that the
full Maurer--Cartan equation is governed by the differential and the
binary bracket alone.
\end{definition}

\begin{theorem}[Pairwise-to-all-point reconstruction]
\ClaimStatusProvedHere
\label{thm:pairwise-to-all-point-reconstruction}
Let $Y$ be in the strict binary regime, and let
\begin{equation}
r(z)\in Y\otimes Y((z^{-1}))
\end{equation}
be a degree-$1$ kernel satisfying the binary Maurer--Cartan equation and
the spectral classical Yang--Baxter equation. For $n\ge 2$ define
\begin{equation}
R_n(z_1,\dots,z_n):=\sum_{1\le i<j\le n} r_{ij}(z_i-z_j).
\end{equation}
Then $R_n$ is Maurer--Cartan for every $n$.
\end{theorem}

\begin{proof}
In the strict binary regime,
\begin{equation}
\operatorname{MC}(R_n)=dR_n+\frac12[R_n,R_n].
\end{equation}
Expanding the bracket, commutators between disjoint pairs vanish because
they act on disjoint tensor factors. The only mixed terms surviving are
those supported on triples $i<j<k$, namely
\begin{equation}
[r_{ij},r_{ik}]+[r_{ij},r_{jk}]+[r_{ik},r_{jk}].
\end{equation}
Thus
\begin{equation}
\operatorname{MC}(R_n)
=
\sum_{i<j}\left(dr_{ij}+\frac12[r_{ij},r_{ij}]\right)
+
\sum_{i<j<k}\Big([r_{ij},r_{ik}]+[r_{ij},r_{jk}]+[r_{ik},r_{jk}]\Big).
\end{equation}
The first sum vanishes by the binary Maurer--Cartan equation and the
second by the classical Yang--Baxter equation. Hence
$\operatorname{MC}(R_n)=0$.
\end{proof}

\begin{remark}[Compact-core consequence]
\label{rem:compact-core-consequence}
Corollary~\ref{cor:one-coefficient-forces-reduced-core} and
Theorem~\ref{thm:pairwise-to-all-point-reconstruction} show that on the
strict polynomial core a single ordered off-diagonal coefficient on the
fundamental line forces the reduced binary kernel and then all higher
ordered corrections. In that regime the compact Yangian core is rigid
rather than mysterious.
\end{remark}

\section{The principal RTT mode tower and weightwise MC4}
\label{sec:typeA-weightwise-mc4}

The next scale is the completion boundary. For
$Y_\hbar(\mathfrak{sl}_N)$ the correct tower is not arbitrary: it is
forced by the mode filtration of the RTT presentation.

\subsection{Principal truncations}

Let
\begin{equation}
T(u)=\mathbf 1+\sum_{r\ge 1}T^{(r)}u^{-r},
\qquad T^{(r)}=(t^{(r)}_{ij})_{1\le i,j\le N},
\end{equation}
with RTT relation
\begin{equation}
R(u-v)T_1(u)T_2(v)=T_2(v)T_1(u)R(u-v),
\qquad R(z)=1-\frac{\hbar P}{z}.
\end{equation}
Define the principal RTT truncation by
\begin{equation}
Y^{\mathrm{RTT}}_{\le M}(\mathfrak{gl}_N)
:=
Y_\hbar(\mathfrak{gl}_N)\Big/\left\langle t^{(r)}_{ij}:r>M\right\rangle.
\end{equation}
Then set
\begin{equation}
Y^{\mathrm{RTT}}_{\le M}(\mathfrak{sl}_N)
:=
Y^{\mathrm{RTT}}_{\le M}(\mathfrak{gl}_N)
\Big/
\left\langle [u^{-r}](\operatorname{qdet}T(u)-1):1\le r\le M\right\rangle.
\end{equation}
The projection dropping the mode $(M+1)$ generators will be denoted
\begin{equation}
p_M:Y^{\mathrm{RTT}}_{\le M+1}(\mathfrak{sl}_N)\twoheadrightarrow
Y^{\mathrm{RTT}}_{\le M}(\mathfrak{sl}_N).
\end{equation}
Assign weight
\begin{equation}
\operatorname{wt}(t^{(r)}_{ij})=r,
\end{equation}
and extend multiplicatively to monomials.

\begin{proposition}[Weight-stability of the algebra tower]
\ClaimStatusProvedHere
\label{prop:weight-stability-algebra-tower}
For every $w\le M$, the map
\begin{equation}
\label{eq:weight-stability-algebra-map}
p_M:
\big(Y^{\mathrm{RTT}}_{\le M+1}(\mathfrak{sl}_N)\big)^{(w)}
\longrightarrow
\big(Y^{\mathrm{RTT}}_{\le M}(\mathfrak{sl}_N)\big)^{(w)}
\end{equation}
is an isomorphism.
\end{proposition}

\begin{proof}
The kernel of $p_M$ is generated by the new mode-$(M+1)$ elements
$t^{(M+1)}_{ij}$ together with the coefficient
$[u^{-(M+1)}](\operatorname{qdet}T(u)-1)$. Every generator of this kernel
has weight $M+1$. Since all generators of the algebra have strictly
positive weight, any element of the kernel has weight at least $M+1$.
Hence the kernel has trivial weight-$w$ part for $w\le M$. Surjectivity
of \eqref{eq:weight-stability-algebra-map} is obvious.
\end{proof}

The preceding proposition immediately explains why the original, global
form of MC4 is too strong for the mode tower. MC4 has since been
resolved via the weightwise approach
(Theorem~\ref{thm:winfty-all-stages-rigidity-closure}).

\begin{proposition}[Failure of unweighted stabilization]
\ClaimStatusProvedHere
\label{prop:failure-unweighted-stabilization}
The inverse system
\begin{equation}
\{\bar B(Y^{\mathrm{RTT}}_{\le M}(\mathfrak{sl}_N))\}_{M\ge 1}
\end{equation}
does not stabilize in ordinary homology degree. In particular,
$H_1\big(\bar B(Y^{\mathrm{RTT}}_{\le M}(\mathfrak{sl}_N))\big)$ does not
stabilize as $M\to\infty$.
\end{proposition}

\begin{proof}
For any augmented algebra $A$,
\begin{equation}
H_1(\bar B(A))\cong I_A/I_A^2,
\end{equation}
where $I_A$ is the augmentation ideal. In
$Y^{\mathrm{RTT}}_{\le M+1}(\mathfrak{sl}_N)$ the classes of the new
mode-$(M+1)$ generators survive in $I/I^2$: the defining RTT relations
have no linear term killing them, and they lie outside $I^2$ by weight
considerations. Under the projection to stage $M$ these new classes
vanish. Hence the maps on $H_1$ are never isomorphisms.
\end{proof}

\begin{remark}[Correct form of MC4: weight filtration on the convolution algebra]
\label{rem:correct-form-of-MC4}
For the principal RTT tower the correct convergence notion is
\emph{weightwise} stabilization, not ordinary stabilization.
In the convolution algebra $\Convstr(C^!_{\mathrm{ch}},\,
P^{\mathrm{ch}})$, the weight grading $\operatorname{wt}(t^{(r)}_{ij})
= r$ induces a descending filtration $F^w$ that is complete,
separated, and respected by all brackets. The MC element $\Theta$
lives in the completed convolution algebra
$\widehat{\Convstr} = \prod_{w \ge 0} F^w / F^{w+1}$, and the
bar-cobar quasi-isomorphism holds weightwise because each weight
piece stabilizes at finite stage (Lemma~\ref{lem:weight-stability-bar-tower}).
This is the form in which MC4 was proved
(Theorem~\ref{thm:winfty-all-stages-rigidity-closure}).
\end{remark}

\subsection{Weightwise stabilization of the bar tower}

Let
\begin{equation}
\bar B_p^{(w)}(M):=
\bar B_p\big(Y^{\mathrm{RTT}}_{\le M}(\mathfrak{sl}_N)\big)^{(w)}
\end{equation}
be the reduced bar complex in bar degree $p$ and total weight $w$.

\begin{lemma}[Weight-stability of the bar tower]
\ClaimStatusProvedHere
\label{lem:weight-stability-bar-tower}
For every $p,w\ge 0$, the transition map
\begin{equation}
\bar B_p^{(w)}(M+1)\longrightarrow \bar B_p^{(w)}(M)
\end{equation}
is an isomorphism for all $M\ge w$.
\end{lemma}

\begin{proof}
A reduced bar tensor $[a_1|\cdots|a_p]$ has total weight
\begin{equation}
\operatorname{wt}(a_1)+\cdots+\operatorname{wt}(a_p)=w.
\end{equation}
Hence each factor has weight at most $w$. If $M\ge w$, then no factor can
involve a mode strictly larger than $M$. By
Proposition~\ref{prop:weight-stability-algebra-tower}, the algebra itself
has stabilized in every weight $\le w$, and therefore the corresponding
bar tensor sector of total weight $w$ has stabilized as well.
\end{proof}

Define the weight-completed bar object by
\begin{equation}
\widehat{\bar B}_{\mathrm{wt}}\big(Y_\hbar(\mathfrak{sl}_N)\big)
:=
\prod_{w\ge 0}\varprojlim_M \bar B^{(w)}\big(Y^{\mathrm{RTT}}_{\le M}(\mathfrak{sl}_N)\big),
\end{equation}
and the weight-completed algebra by
\begin{equation}
\widehat Y^{\mathrm{RTT}}_{\mathrm{wt}}(\mathfrak{sl}_N)
:=
\prod_{w\ge 0}Y_\hbar(\mathfrak{sl}_N)^{(w)}.
\end{equation}

\begin{theorem}[Weightwise MC4 for the principal RTT tower]
\ClaimStatusProvedHere
\label{thm:weightwise-MC4-principal-RTT}
Assume that each finite stage
$Y^{\mathrm{RTT}}_{\le M}(\mathfrak{sl}_N)$ lies in the finite-stage
bar--cobar regime of the monograph. Then:
\begin{enumerate}
 \item $\widehat{\bar B}_{\mathrm{wt}}(Y_\hbar(\mathfrak{sl}_N))$ is a
 separated complete conilpotent dg coalgebra.
 \item The completed bar differential is continuous for the product
 topology.
 \item The completed cobar construction satisfies a quasi-isomorphism
 \begin{equation}
 \widehat\Omega\Big(
 \widehat{\bar B}_{\mathrm{wt}}(Y_\hbar(\mathfrak{sl}_N))
 \Big)
 \xrightarrow{\sim}
 \widehat Y^{\mathrm{RTT}}_{\mathrm{wt}}(\mathfrak{sl}_N).
 \end{equation}
\end{enumerate}
\end{theorem}

\begin{proof}
By Lemma~\ref{lem:weight-stability-bar-tower}, the inverse system in each
fixed weight $w$ is eventually constant. Hence each weight piece of the
limit exists and is canonically identified with the stable value at any
stage $M\ge w$. Taking the product over all weights gives a separated,
complete graded coalgebra.

The bar differential preserves total weight, so it acts independently on
each stable weight piece and is therefore continuous for the product
topology.

For each fixed weight $w$, the bar--cobar quasi-isomorphism stabilizes at
finite stage $M\ge w$ by the same reasoning: every expression of total
weight $w$ only uses generators of weight at most $w$. Thus in each
weight piece the completed cobar of the completed bar is canonically
identified with the stable finite-stage cobar, which is quasi-isomorphic
to the stable finite-stage algebra. Taking the product over $w$ yields
the stated quasi-isomorphism.
\end{proof}

The next statement is a recognition theorem: it identifies the
completed bar dual as a dg-shifted Yangian, i.e., the binary collision
residue of the universal MC element. In the language of the holographic
chapter (\S\ref{subsec:holographic-modular-koszul-datum}), the
MC kernel $r(z)$ appearing below is
$\operatorname{Res}^{\mathrm{coll}}_{0,2}(\Theta_\cA)$, the same
object that seeds the full shadow obstruction tower.

\begin{theorem}[Recognition of the completed dg-shifted Yangian]
\ClaimStatusProvedHere
\label{thm:recognition-completed-dg-shifted-yangian}
Assume in addition that the principal RTT tower carries a continuous
derivation implementing spectral translation and compatible with all
finite-stage truncations. Let $A^!$ be the $A_\infty$-bar dual of
$\widehat Y^{\mathrm{RTT}}_{\mathrm{wt}}(\mathfrak{sl}_N)$. Then $A^!$
admits a canonical dg-shifted Yangian structure in the sense of
\cite{DNP25}: the translated universal twisting cochain determines a
Maurer--Cartan kernel
\begin{equation}
r(z)\in A^!\widehat\otimes A^!((z^{-1})),
\end{equation}
and the deconcatenation coproduct on the completed bar object induces the
twisted coproduct of the dg-shifted Yangian.
\end{theorem}

\begin{proof}
Let
\begin{equation}
\tau:\widehat{\bar B}_{\mathrm{wt}}(Y_\hbar(\mathfrak{sl}_N))\to A^!
\end{equation}
be the universal twisting cochain. By construction it satisfies the
Maurer--Cartan equation in the convolution dg Lie algebra,
\begin{equation}
d\tau+\tau\star\tau=0.
\end{equation}
The continuous spectral derivation on the completed Yangian induces a
coderivation on the completed bar coalgebra and therefore an
$A_\infty$-derivation on $A^!$. Exponentiating gives a translation family
$\tau_z$. Its bar-degree-$2$ component is the desired kernel $r(z)$. The
usual pushforward of deconcatenation along the twisting cochain produces a
twisted coproduct. Coassociativity and the Maurer--Cartan identities are
immediate from the coalgebra identities and from the twisting-cochain
Maurer--Cartan equation.
\end{proof}

\section{Derived KR telescopes and the noncompact packet}
\label{sec:typeA-derived-KR-telescopes}

The compact polynomial core is not the whole category. The genuinely
noncompact packet is generated by asymptotic and prefundamental objects,
and in type~$A$ both arise from the same Kirillov--Reshetikhin chain.

Fix a simple node $i\in\{1,\dots,N-1\}$ and write
\begin{equation}
W_l:=W^{(i)}_{l,0}
\end{equation}
for the Kirillov--Reshetikhin module of level $l$ and zero shift. Zhang
constructs an inductive system
\begin{equation}
W_0\xhookrightarrow{F_{1,0}}W_1\xhookrightarrow{F_{2,1}}W_2\xhookrightarrow{}\cdots
\end{equation}
and shows that the Yangian action is affine-linear in the stage parameter
along this system~\cite{Zhang24}. The new observation is that this chain
already carries its own derived compactification.

\subsection{Ordinary asymptotic telescopes}

\begin{definition}[Ordinary asymptotic telescope]
\label{def:ordinary-asymptotic-telescope}
Assume that for every Yangian generator $s$ there exist linear maps
\begin{equation}
A_s^l,B_s^l:W_l\longrightarrow W_{l+1}
\end{equation}
such that for all sufficiently large $k>l$,
\begin{equation}
\label{eq:ordinary-asymptotic-stage-action}
s\,F_{k,l}=F_{k,l+1}(A_s^l+\lambda B_s^l)\big|_{\lambda=k}.
\end{equation}
Define the \emph{ordinary asymptotic telescope complex}
\begin{equation}
\label{eq:ordinary-telescope-definition}
\mathrm{Tel}^{\mathrm{ord}}_{i,\lambda}
:=
\operatorname{Cone}\!\left(
\bigoplus_{l\ge 0}W_l
\xrightarrow{\ \mathrm{id}-F\ }
\bigoplus_{l\ge 0}W_l
\right),
\end{equation}
where $F$ is induced by the structure maps $F_{l+1,l}$, and the action of
$s$ on the $l$-th summand is declared to be $A_s^l+\lambda B_s^l$.
\end{definition}

\begin{theorem}[Derived realization of ordinary asymptotic modules]
\ClaimStatusProvedHere
\label{thm:derived-realization-ordinary-asymptotic}
With the notation above, $\mathrm{Tel}^{\mathrm{ord}}_{i,\lambda}$ is a
complex of $Y_\hbar(\mathfrak{sl}_N)$-modules, and
\begin{equation}
H^{-1}(\mathrm{Tel}^{\mathrm{ord}}_{i,\lambda})=0,
\qquad
H^0(\mathrm{Tel}^{\mathrm{ord}}_{i,\lambda})
\cong
L\!\left(\frac{\Psi_{i,y}}{\Psi_{i,x}}\right)
\end{equation}
whenever $\lambda=(y-x)/\hbar$ and the spectral shift is normalized so
that the asymptotic parameter is $y-x$.
\end{theorem}

\begin{proof}
The inductive-system compatibility of the maps $A_s^l$ and $B_s^l$ implies
that every generator action commutes with the telescope differential
$\mathrm{id}-F$. Hence the complex carries a Yangian action.

Because each transition map $F_{l+1,l}$ is injective, the standard
cohomology calculation for a mapping telescope gives
$H^{-1}=0$ and identifies $H^0$ with the inductive limit vector space
$\varinjlim_l W_l$. The action of a generator on this inductive limit is,
by definition, the asymptotic action obtained by specializing the affine
linear stage formula \eqref{eq:ordinary-asymptotic-stage-action} at the
parameter $\lambda$. This is Zhang's asymptotic module
construction~\cite{Zhang24}.
\end{proof}

\subsection{Boundary telescopes and negative prefundamentals}

The shifted chart is different in appearance but not in origin: it is the
same KR chain after Cartan renormalization.

\begin{definition}[Boundary telescope]
\label{def:boundary-telescope}
Assume that for the Cartan-doubled Yangian $Y_\infty(\mathfrak{sl}_N)$
there are linear maps
\begin{equation}
C_t^l,D_t^l:W_l\longrightarrow W_{l+1}
\end{equation}
for every generator $t$ such that for sufficiently large $k>l$,
\begin{equation}
\label{eq:boundary-stage-action}
t\,F_{k,l}=F_{k,l+1}(C_t^l+qD_t^l)\big|_{q=k^{-1}}.
\end{equation}
Define the \emph{boundary telescope complex}
\begin{equation}
\mathrm{Tel}^{\mathrm{bdry}}_{i,q}
:=
\operatorname{Cone}\!\left(
\bigoplus_{l\ge 0}W_l
\xrightarrow{\ \mathrm{id}-F\ }
\bigoplus_{l\ge 0}W_l
\right),
\end{equation}
with $Y_\infty(\mathfrak{sl}_N)$-action on the $l$-th summand given by
$C_t^l+qD_t^l$.
\end{definition}

\begin{theorem}[Derived realization of negative prefundamentals]
\ClaimStatusProvedHere
\label{thm:derived-realization-negative-prefundamental}
The complex $\mathrm{Tel}^{\mathrm{bdry}}_{i,q}$ is a complex of
$Y_\infty(\mathfrak{sl}_N)$-modules. At the boundary point $q=0$ one has
\begin{equation}
H^{-1}(\mathrm{Tel}^{\mathrm{bdry}}_{i,0})=0,
\qquad
H^0(\mathrm{Tel}^{\mathrm{bdry}}_{i,0})\cong L^-_{i,0},
\end{equation}
and after spectral shift,
\begin{equation}
H^0(\tau_x^*\mathrm{Tel}^{\mathrm{bdry}}_{i,0})\cong L^-_{i,x}.
\end{equation}
\end{theorem}

\begin{proof}
As in Theorem~\ref{thm:derived-realization-ordinary-asymptotic}, the
compatibility of the coefficient maps with the inductive system shows that
each generator acts by a chain map. Injectivity of the transition maps
again gives $H^{-1}=0$ and identifies $H^0$ with the inductive limit
vector space.

The $q=0$ specialization of the generator action is exactly the renormalized
boundary action $t\mapsto C_t$. By Zhang's construction this action
factors through the anti-dominantly shifted Yangian and identifies the
inductive limit with the negative prefundamental module~\cite{Zhang24}.
\end{proof}

\begin{definition}[Baxter packet]
\label{def:baxter-packet}
Let $\widehat{\mathcal O}$ be a completed or coderived enhancement of the
Yangian category. The \emph{Baxter packet} is the localizing tensor
subcategory
\begin{equation}
\mathcal G_{\mathrm{Bax}}
:=
\operatorname{Loc}^{\otimes}
\Big(
\mathcal O^{\omega}_{\mathrm{poly}}
\cup
\{L^-_{i,a}\}_{1\le i\le N-1,\,a\in\mathbb C}
\Big).
\end{equation}
\end{definition}

\begin{theorem}[Baxter-envelope criterion]
\ClaimStatusProvedHere
\label{thm:baxter-envelope-criterion}
Assume:
\begin{enumerate}
 \item every asymptotic module is realized by an ordinary asymptotic
 telescope built from compact KR modules;
 \item every standard object of every truncated shifted Yangian admits
 a finite filtration with subquotients given by tensor products of
 compact polynomial objects and negative prefundamentals;
 \item every simple object of $\widehat{\mathcal O}$ is a quotient,
 extension, or filtered homotopy colimit of such standard and
 asymptotic modules;
 \item $\widehat{\mathcal O}$ is generated from its heart under cones,
 shifts, and filtered colimits.
\end{enumerate}
Then
\begin{equation}
\widehat{\mathcal O}=\mathcal G_{\mathrm{Bax}}.
\end{equation}
\end{theorem}

\begin{proof}
By Theorem~\ref{thm:derived-realization-ordinary-asymptotic}, every
asymptotic module belongs to the localizing tensor closure of the compact
KR packet, hence to $\mathcal G_{\mathrm{Bax}}$. By hypothesis~(2), the
same holds for all standard truncated objects. Hypothesis~(3) expresses
every simple object in terms of these building blocks, and
hypothesis~(4) then implies that the whole stable category lies in the
localizing tensor closure generated by them. The reverse inclusion is
built into the definition of $\mathcal G_{\mathrm{Bax}}$.
\end{proof}

\section{The Baxter--Rees family and its formal neighborhood}
\label{sec:typeA-baxter-rees-family}

The two telescopes constructed above do not merely meet at a point. They
assemble into a single affine family carrying the generic asymptotic chart
and the boundary prefundamental chart at once.

\begin{definition}[Baxter--Rees family]
\label{def:baxter-rees-family}
Let
\begin{equation}
W_\infty:=\varinjlim_l W_l.
\end{equation}
Define the free $\mathbb C[q]$-module
\begin{equation}
\mathcal E_i:=W_\infty\otimes_{\mathbb C}\mathbb C[q].
\end{equation}
For every generator $t$ of $Y_\infty(\mathfrak{sl}_N)$ set
\begin{equation}
\rho_q(t):=C_t+qD_t\in \operatorname{End}_{\mathbb C[q]}(\mathcal E_i).
\end{equation}
\end{definition}

The point of this definition is that it is not merely formal. The stage
formula \eqref{eq:boundary-stage-action} determines a genuine algebraic
family.

\begin{theorem}[Algebraicity of the Baxter--Rees family]
\ClaimStatusProvedHere
\label{thm:algebraicity-baxter-rees-family}
The assignment $t\mapsto \rho_q(t)$ extends to a
$\mathbb C[q]$-linear representation
\begin{equation}
\rho_q:Y_\infty(\mathfrak{sl}_N)
\longrightarrow
\operatorname{End}_{\mathbb C[q]}(\mathcal E_i).
\end{equation}
Hence $\mathcal E_i$ is a flat algebraic family of
$Y_\infty(\mathfrak{sl}_N)$-modules over $\mathbb A^1_q$.
\end{theorem}

\begin{proof}
Take any defining polynomial relation
$\mathscr R(t_1,\dots,t_m)=0$ of the algebra. Restrict the operator
obtained by substituting $t_a\mapsto C_{t_a}^l+qD_{t_a}^l$ to the image of
$W_l$. Because each generator contributes at most one power of $q$, the
result is a polynomial in $q$ with values in a finite-dimensional vector
space:
\begin{equation}
\mathscr R_l(q):W_l\longrightarrow W_{l+m}.
\end{equation}
For every sufficiently large integer $k>l$ we may evaluate at $q=k^{-1}$,
and by the stage formula \eqref{eq:boundary-stage-action} the resulting
operator is exactly the action of the genuine finite-stage module obtained
via the renormalized chart. Hence $\mathscr R_l(k^{-1})=0$ for infinitely
many values of $k$. A polynomial over $\mathbb C$ vanishing on an
infinite set is zero, so $\mathscr R_l(q)\equiv 0$ for every $l$. This
proves the algebra relations. Flatness is immediate from the fact that
$\mathcal E_i$ is free over $\mathbb C[q]$.
\end{proof}

\begin{theorem}[Generic and boundary fibers]
\ClaimStatusProvedHere
\label{thm:generic-and-boundary-fibers}
For every $a\in \mathbb C^\times$ the fiber
\begin{equation}
(\mathcal E_i)_a:=\mathcal E_i\otimes_{\mathbb C[q]}\mathbb C_{q=a}
\end{equation}
is canonically isomorphic to the pullback of the ordinary asymptotic
module along the Cartan renormalization map,
\begin{equation}
(\mathcal E_i)_a
\cong
\theta_{i,a^{-1}}^*L\!\left(\frac{\Psi_{i,a^{-1}}}{\Psi_{i,0}}\right).
\end{equation}
At the boundary,
\begin{equation}
(\mathcal E_i)_0\cong L^-_{i,0}.
\end{equation}
After spectral shift,
\begin{equation}
\tau_x^*(\mathcal E_i)_a
\cong
\theta_{i,a^{-1}}^*L\!\left(\frac{\Psi_{i,x+a^{-1}}}{\Psi_{i,x}}\right),
\qquad
\tau_x^*(\mathcal E_i)_0\cong L^-_{i,x}.
\end{equation}
\end{theorem}

\begin{proof}
The generic fiber action is obtained by specializing
$\rho_q(t)=C_t+qD_t$ at $q=a$. By construction of the renormalized chart,
this specialization is the same as the pullback of the ordinary
asymptotic action at parameter $a^{-1}$ along the Cartan renormalization
map $\theta_{i,a^{-1}}$. The $q=0$ statement is exactly the boundary case
of Theorem~\ref{thm:derived-realization-negative-prefundamental}. The
spectral-shift statements follow from functoriality of the translation
automorphisms.
\end{proof}

The family itself also has a derived realization.

\begin{theorem}[Derived realization of the Baxter--Rees family]
\ClaimStatusProvedHere
\label{thm:derived-realization-baxter-rees-family}
Let
\begin{equation}
\mathfrak T_i^{\mathrm{Bax}}
:=
\operatorname{Cone}\!\left(
\bigoplus_{l\ge 0}W_l\otimes \mathbb C[q]
\xrightarrow{\ \mathrm{id}-F\ }
\bigoplus_{l\ge 0}W_l\otimes \mathbb C[q]
\right),
\end{equation}
where the action of $t\in Y_\infty(\mathfrak{sl}_N)$ on the $l$-th summand
is $C_t^l+qD_t^l$. Then
$\mathfrak T_i^{\mathrm{Bax}}$ is a complex of
$Y_\infty(\mathfrak{sl}_N)\otimes\mathbb C[q]$-modules, and
\begin{equation}
H^{-1}(\mathfrak T_i^{\mathrm{Bax}})=0,
\qquad
H^0(\mathfrak T_i^{\mathrm{Bax}})\cong \mathcal E_i.
\end{equation}
\end{theorem}

\begin{proof}
The proof is the same telescope computation as before. Compatibility of
$C_t^l$ and $D_t^l$ with the inductive system implies that every generator
acts by a chain map, so $\mathfrak T_i^{\mathrm{Bax}}$ is a module complex.
Injectivity of the transition maps gives vanishing of $H^{-1}$ and
identifies $H^0$ with
\begin{equation}
\varinjlim_l (W_l\otimes \mathbb C[q])=W_\infty\otimes \mathbb C[q]=\mathcal E_i.
\end{equation}
The induced action is precisely $t\mapsto C_t+qD_t$.
\end{proof}

\begin{corollary}[Formal neighborhood generated by the compact KR packet]
\ClaimStatusProvedHere
\label{cor:formal-neighborhood-generated-by-KR}
For every $n\ge 0$ the jet module $\mathcal E_i/(q^{n+1})$ belongs to the
localizing subcategory generated by the compact KR packet
$\{W_l\}_{l\ge 0}$. Consequently the full formal neighborhood
\begin{equation}
\widehat{\mathcal E}_i:=\varprojlim_n \mathcal E_i/(q^{n+1})
\end{equation}
lies in the $q$-adically completed localizing category generated by the
compact KR packet.
\end{corollary}

\begin{proof}
The complex $\mathfrak T_i^{\mathrm{Bax}}$ is built from direct sums of
$W_l\otimes\mathbb C[q]$, so it belongs to the localizing subcategory
generated by the compact KR packet. Passing to the quotient modulo
$q^{n+1}$ preserves this property, and taking $H^0$ yields
$\mathcal E_i/(q^{n+1})$. The final statement follows by inverse limit.
\end{proof}

The family comes with a canonical first-order boundary class.

\begin{definition}[Baxter--Kodaira--Spencer class]
\label{def:baxter-kodaira-spencer-class}
Let
\begin{equation}
M_i^-:=(\mathcal E_i)_0\cong L^-_{i,0}.
\end{equation}
The reduction of $\mathcal E_i$ modulo $q^2$ gives an exact sequence
\begin{equation}
0\longrightarrow q\mathcal E_i/q^2\mathcal E_i
\longrightarrow \mathcal E_i/q^2\mathcal E_i
\longrightarrow \mathcal E_i/q\mathcal E_i
\longrightarrow 0.
\end{equation}
Multiplication by $q$ identifies $q\mathcal E_i/q^2\mathcal E_i$ with
$M_i^-$. The resulting self-extension class in
$\operatorname{Ext}^1_{Y_\infty}(M_i^-,M_i^-)$ is called the
\emph{Baxter--Kodaira--Spencer class} and is denoted $\kappa_i$.
\end{definition}

\begin{theorem}[Concrete cocycle representative]
\ClaimStatusProvedHere
\label{thm:concrete-cocycle-representative}
The class $\kappa_i$ is represented by the $1$-cocycle
\begin{equation}
\delta_i(t):=D_t
\end{equation}
relative to the boundary action $\rho_0(t)=C_t$. Equivalently,
\begin{equation}
\rho_q(t)=\rho_0(t)+q\,\delta_i(t)\pmod{q^2}.
\end{equation}
\end{theorem}

\begin{proof}
Modulo $q^2$ the family has the form
$\rho_q=\rho_0+q\delta_i$. Since $\rho_q$ is an algebra representation,
\begin{equation}
\rho_q(ab)=\rho_q(a)\rho_q(b)\pmod{q^2}.
\end{equation}
Comparing coefficients of $q$ yields
\begin{equation}
\delta_i(ab)=\rho_0(a)\delta_i(b)+\delta_i(a)\rho_0(b),
\end{equation}
which is exactly the cocycle condition describing the associated
square-zero self-extension.
\end{proof}

\section{Weightwise continuation and boundary tensor geometry}
\label{sec:typeA-boundary-tensor-geometry}

The Baxter--Rees family improves the object-theoretic picture, but the
factorization problem lives one level higher: the boundary family must
also carry braidings and associators. This is where the polynomial
$R$-matrix theory of shifted Yangians and Zhang's Theta-associators enter
\cite{HZ24,ZhangTheta24}. The first step is a purely algebraic
continuation principle.

\begin{lemma}[Weightwise polynomial continuation]
\ClaimStatusProvedHere
\label{lem:weightwise-polynomial-continuation}
Let $M$ and $N$ be free $\mathbb C[q]$-modules equipped with direct-sum
decompositions into finite-dimensional weight spaces,
\begin{equation}
M=\bigoplus_{\gamma} M_\gamma,
\qquad N=\bigoplus_{\gamma} N_\gamma.
\end{equation}
Suppose that for each weight $\gamma$ one is given a map on the punctured
chart
\begin{equation}
F_\gamma(q):M_\gamma\otimes \mathbb C[q,q^{-1}]
\longrightarrow N_\gamma\otimes \mathbb C[q,q^{-1}]
\end{equation}
such that, in some bases of the finite-dimensional spaces $M_\gamma$ and
$N_\gamma$, all matrix coefficients are polynomials in $q$. Then:
\begin{enumerate}
 \item there exists a unique $\mathbb C[q]$-linear map
 \begin{equation}
 F:M\to N
 \end{equation}
 whose restriction to each weight space is the polynomial extension of
 $F_\gamma(q)$;
 \item if a finite polynomial identity in such maps holds on a
 Zariski-dense subset of $\mathbb A^1_q$, then it holds identically on
 all of $\mathbb A^1_q$.
\end{enumerate}
\end{lemma}

\begin{proof}
Because every weight space is finite-dimensional, the hypothesis gives a
matrix with polynomial coefficients on each weight space. Evaluating
these polynomials defines the unique extension over $\mathbb C[q]$.
Uniqueness is immediate because a polynomial function is determined by its
values on a Zariski-dense subset. Any finite algebraic identity among
such maps is, weightwise, an identity of matrices with polynomial
coefficients; if it holds on a dense subset, every coefficient polynomial
vanishes identically.
\end{proof}

\subsection{Polynomial continuation of $R$-matrices}

The next theorem isolates the precise input required from the literature.
It applies whenever the normalized intertwiners on the punctured chart are
known weightwise to be polynomial in the boundary coordinate $q$.

\begin{theorem}[Continuation of polynomial $R$-matrices to the boundary]
\ClaimStatusProvedHere
\label{thm:continuation-of-polynomial-R-matrices}
Fix a compact polynomial module $M\in \mathcal O^{\omega}_{\mathrm{poly}}$.
Assume that on the punctured chart $q\neq 0$ there exists a normalized
intertwiner
\begin{equation}
\mathscr R_{M,i}(u,q):
M\otimes (\mathcal E_i)_q
\longrightarrow
(\mathcal E_i)_q\otimes M
\end{equation}
such that:
\begin{enumerate}
 \item on each finite-dimensional weight space, the matrix
 coefficients of $\mathscr R_{M,i}(u,q)$ are polynomials in $q$;
 \item on the punctured chart the family satisfies the intertwining
 equation and the Yang--Baxter equation with the compact polynomial
 packet.
\end{enumerate}
Then:
\begin{enumerate}
 \item $\mathscr R_{M,i}(u,q)$ extends uniquely to a
 $\mathbb C[q]$-linear family on all of $\mathbb A^1_q$;
 \item its specialization at $q=0$ defines a boundary intertwiner
 \begin{equation}
 \mathscr R_{M,i}^{\partial}(u):
 M\otimes L^-_{i,0}
 \longrightarrow
 L^-_{i,0}\otimes M;
 \end{equation}
 \item the extended family satisfies the intertwining and
 Yang--Baxter equations for every $q$, in particular at the boundary.
\end{enumerate}
\end{theorem}

\begin{proof}
Apply Lemma~\ref{lem:weightwise-polynomial-continuation} to the weightwise
matrix coefficients of $\mathscr R_{M,i}(u,q)$. This yields the unique
extension over $\mathbb C[q]$. The specialized map at $q=0$ is then
well-defined and acts on the boundary fiber
$(\mathcal E_i)_0\cong L^-_{i,0}$ by
Theorem~\ref{thm:generic-and-boundary-fibers}. The intertwining and
Yang--Baxter identities are polynomial identities in the entries of the
family, hence they extend from the punctured chart to all of
$\mathbb A^1_q$ by the second part of
Lemma~\ref{lem:weightwise-polynomial-continuation}.
\end{proof}

\begin{remark}[Why the hypotheses are natural]
\label{rem:why-R-matrix-hypotheses-are-natural}
For shifted Yangians, Hernandez--Zhang prove existence and uniqueness of
normalized intertwiners for tensor products involving positive or negative
prefundamental modules and prove polynomiality in the spectral
parameter~\cite{HZ24}. Zhang's Theta-series formalism then supplies the
polynomiality mechanism on the shifted chart and the decomposition of
intertwiners through one-dimensional modules~\cite{ZhangTheta24}. The
present theorem isolates the exact algebraic consequence of that
polynomiality: once the weightwise regularity is known, the boundary
specialization is automatic.
\end{remark}

\subsection{Continuation of Theta-associators}

The same continuation principle applies one categorical level higher.

\begin{theorem}[Continuation of Theta-associators to the boundary]
\ClaimStatusProvedHere
\label{thm:continuation-of-theta-associators}
Fix compact polynomial modules $L,M\in \mathcal O^{\omega}_{\mathrm{poly}}$.
Assume that on the punctured chart $q\neq 0$ there exists a family of
isomorphisms
\begin{equation}
\Theta_{L,M,i}(q):
(L\otimes M)\otimes (\mathcal E_i)_q
\longrightarrow
L\otimes (M\otimes (\mathcal E_i)_q)
\end{equation}
such that:
\begin{enumerate}
 \item on each weight component its matrix coefficients are
 polynomials in $q$;
 \item on the punctured chart the family satisfies the associator
 identities, naturality, and the mixed hexagon compatibilities with
 the continued $R$-matrices.
\end{enumerate}
Then:
\begin{enumerate}
 \item $\Theta_{L,M,i}(q)$ extends uniquely to a $\mathbb C[q]$-linear
 family on all of $\mathbb A^1_q$;
 \item its specialization at the boundary defines a boundary
 associator
 \begin{equation}
 \Theta^{\partial}_{L,M,i}:
 (L\otimes M)\otimes L^-_{i,0}
 \longrightarrow
 L\otimes (M\otimes L^-_{i,0});
 \end{equation}
 \item the extended family satisfies the associativity and mixed
 compatibility identities for every $q$, in particular at the
 boundary.
\end{enumerate}
\end{theorem}

\begin{proof}
The proof is identical to that of
Theorem~\ref{thm:continuation-of-polynomial-R-matrices}. Every relevant
identity is a finite polynomial relation among the weightwise matrix
coefficients of the family, so the continuation lemma applies verbatim.
\end{proof}

\begin{corollary}[Formal braided boundary germ]
\ClaimStatusProvedHere
\label{cor:formal-braided-boundary-germ}
For each simple node $i$, the boundary fiber $L^-_{i,0}$ together with the
specialized intertwiners and associators
\begin{equation}
\mathscr R_{M,i}^{\partial}(u),
\qquad
\Theta^{\partial}_{L,M,i}
\end{equation}
defines a braided monoidal boundary object inside the formal neighborhood
\begin{equation}
\widehat{\mathcal E}_i=\varprojlim_n \mathcal E_i/(q^{n+1}).
\end{equation}
More precisely, the completed localizing envelope generated by the KR
packet acquires a distinguished formal boundary germ at each simple root.
\end{corollary}

\begin{proof}
The specialized $R$-matrices and associators satisfy the required
identities by
Theorems~\ref{thm:continuation-of-polynomial-R-matrices}
and~\ref{thm:continuation-of-theta-associators}. The formal neighborhood
statement is exactly Corollary~\ref{cor:formal-neighborhood-generated-by-KR}.
\end{proof}

\subsection{Linearized boundary tensor constraints}

The continuation theorems say more than existence: they endow the boundary
with a tangent monoidal structure.

\begin{definition}[Boundary tangent tensors]
\label{def:boundary-tangent-tensors}
Assume the families of
Theorems~\ref{thm:continuation-of-polynomial-R-matrices}
and~\ref{thm:continuation-of-theta-associators} are defined. Set
\begin{equation}
\dot{\mathscr R}_{M,i}(u)
:=
\left.\frac{d}{dq}\right|_{q=0}\mathscr R_{M,i}(u,q),
\qquad
\dot\Theta_{L,M,i}
:=
\left.\frac{d}{dq}\right|_{q=0}\Theta_{L,M,i}(q).
\end{equation}
\end{definition}

\begin{theorem}[Linearized boundary equations]
\ClaimStatusProvedHere
\label{thm:linearized-boundary-equations}
The first-order tensors
$\dot{\mathscr R}_{M,i}(u)$ and $\dot\Theta_{L,M,i}$ satisfy the
linearized intertwining, Yang--Baxter, and pentagon equations around the
boundary structure.

More explicitly:
\begin{enumerate}
 \item if $\mathscr R_{12}(u,q)$,
 $\mathscr R_{13}(u+v,q)$, and $\mathscr R_{23}(v,q)$ satisfy the
 Yang--Baxter equation for all $q$, then at $q=0$ one has
 \begin{align}
 &\dot{\mathscr R}_{12}(u)
 \mathscr R^{\partial}_{13}(u+v)
 \mathscr R^{\partial}_{23}(v)
 +
 \mathscr R^{\partial}_{12}(u)
 \dot{\mathscr R}_{13}(u+v)
 \mathscr R^{\partial}_{23}(v)
 +
 \mathscr R^{\partial}_{12}(u)
 \mathscr R^{\partial}_{13}(u+v)
 \dot{\mathscr R}_{23}(v)
 \nonumber\\
 &=
 \dot{\mathscr R}_{23}(v)
 \mathscr R^{\partial}_{13}(u+v)
 \mathscr R^{\partial}_{12}(u)
 +
 \mathscr R^{\partial}_{23}(v)
 \dot{\mathscr R}_{13}(u+v)
 \mathscr R^{\partial}_{12}(u)
 +
 \mathscr R^{\partial}_{23}(v)
 \mathscr R^{\partial}_{13}(u+v)
 \dot{\mathscr R}_{12}(u);
 \label{eq:linearized-YB-boundary}
 \end{align}
 \item the derivative $\dot\Theta$ is a cocycle for the linearization
 of the pentagon relation around the specialized associator
 $\Theta^{\partial}$; equivalently, the first-order term of the
 pentagon identity vanishes.
\end{enumerate}
\end{theorem}

\begin{proof}
Differentiate the corresponding identities of the extended families with
respect to $q$ and evaluate at $q=0$. Since differentiation is a
derivation, the first-order term is obtained by replacing one factor at a
time by its derivative and then summing over all such choices. This gives
\eqref{eq:linearized-YB-boundary} and the linearized pentagon equation.
The same argument applies to the intertwining equation.
\end{proof}

\begin{remark}[Boundary deformation complex]
\label{rem:boundary-deformation-complex}
The pair consisting of the Baxter--Kodaira--Spencer class $\kappa_i$ and
the tangent monoidal tensors $\dot{\mathscr R}_{M,i}$,
$\dot\Theta_{L,M,i}$ should be regarded as a single class in the tangent
complex governing braided monoidal deformations of the boundary object
$L^-_{i,0}$. The present appendix does not develop that full deformation
complex, but it provides the concrete cocycles from which it must be
built.
\end{remark}

\section{Compact-to-completed extension and the type-$A$ DK reduction}
\label{sec:typeA-DK-reduction}

The previous sections produce the geometry and tensor data of the
boundary. The remaining step is categorical: once the compact core is
matched with its factorization or spectral counterpart, the Baxter--Rees
package should extend that equivalence to the completed/coderived ambient.

\begin{theorem}[Compact-to-completed extension principle]
\ClaimStatusProvedHere
\label{thm:compact-to-completed-extension-principle}
Let
\begin{equation}
F^\omega:\mathcal C^\omega\xrightarrow{\sim}\mathcal D^\omega
\end{equation}
be an exact braided monoidal equivalence of small idempotent-complete
stable categories. Assume:
\begin{enumerate}
 \item $\widehat{\mathcal C}$ and $\widehat{\mathcal D}$ are compactly
 generated presentable stable categories with compact objects
 $\mathcal C^\omega$ and $\mathcal D^\omega$;
 \item each is obtained from the Ind-completion of its compact objects
 by localizing at a corresponding tensor ideal of coacyclic or
 completion objects;
 \item $\operatorname{Ind}(F^\omega)$ identifies those tensor ideals.
\end{enumerate}
Then $F^\omega$ extends uniquely to a colimit-preserving exact braided
monoidal equivalence
\begin{equation}
\widehat F:\widehat{\mathcal C}\xrightarrow{\sim}\widehat{\mathcal D}.
\end{equation}
\end{theorem}

\begin{proof}
Ind-completion turns an exact equivalence of compact stable categories
into a colimit-preserving exact equivalence of compactly generated stable
categories. Since the localizing tensor ideals correspond under this
Ind-extension, the universal property of Verdier localization gives the
descended equivalence on the quotients. Uniqueness is automatic because
the compact objects generate.
\end{proof}

The resulting reduction theorem is the precise type-$A$ version of the
thread's main conclusion.

\begin{theorem}[Type-$A$ reduction of the full derived DK frontier]
\ClaimStatusProvedHere
\label{thm:typeA-reduction-full-derived-DK-frontier}
Assume the following input.
\begin{enumerate}
 \item \textbf{Compact-core equivalence.}
 There is an exact braided monoidal equivalence on the compact
 polynomial core, compatible with spectral shifts and the ordered
 vector seed, between the Yangian side and a factorization or spectral
 quantum-group side. In particular, it matches the KR packet.
 \item \textbf{Weightwise MC4.}
 The principal RTT tower is used on the Yangian side, so that
 Theorem~\ref{thm:weightwise-MC4-principal-RTT} applies and the
 completed dg-shifted Yangian exists by
 Theorem~\ref{thm:recognition-completed-dg-shifted-yangian}.
 \item \textbf{Baxter envelope.}
 The completed or coderived Yangian category satisfies the hypotheses
 of Theorem~\ref{thm:baxter-envelope-criterion}.
 \item \textbf{Boundary tensor regularity.}
 The polynomiality hypotheses of
 Theorems~\ref{thm:continuation-of-polynomial-R-matrices}
 and~\ref{thm:continuation-of-theta-associators} hold on the punctured
 boundary chart.
 \item \textbf{Matching localizing ideals.}
 The completion or coacyclic tensor ideals match under the compact-core
 equivalence.
\end{enumerate}
Then the compact-core equivalence extends to an exact braided monoidal
equivalence on the completed/coderived Baxter envelope. Under this
extension, the KR telescopes, the Baxter--Rees family, the negative
prefundamental objects, the formal neighborhoods, the continued
$R$-matrices, the continued Theta-associators, and the boundary tangent
classes are all transported functorially.
\end{theorem}

\begin{proof}
Theorem~\ref{thm:baxter-envelope-criterion} identifies the full completed
Yangian side with the localizing tensor envelope of the compact KR and
negative prefundamental packet. Theorem~\ref{thm:compact-to-completed-extension-principle}
then extends the compact-core equivalence to the completed/coderived
ambient. Exactness guarantees that mapping cones are preserved, so the KR
telescope constructions and hence the Baxter--Rees family are transported.
The continuation of $R$-matrices and Theta-associators is functorial by
construction, and the boundary tangent classes are obtained by first-order
reduction of transported families.
\end{proof}

\begin{corollary}[What remains after the present appendix]
\ClaimStatusProvedHere
\label{cor:what-remains-after-present-appendix}
After the constructions and reductions above, the full type-$A$ frontier
is no longer ``extend from evaluation modules to all of category
$\mathcal O$''. It reduces to four concrete verification problems:
\begin{enumerate}
 \item verify the compact-core equivalence on the KR packet;
 \item verify the translation-compatible finite-stage bar--cobar
 package on the chosen RTT truncations;
 \item verify the Baxter-envelope hypotheses on the completed/coderived
 Yangian category;
 \item verify the weightwise polynomial regularity of the boundary
 $R$-matrix and Theta-associator families on the punctured chart.
\end{enumerate}
Once these are known, the extension to the full completed/coderived
factorization category is formal.
\end{corollary}

\section{Platonic synthesis and modular outlook}
\label{sec:typeA-platonic-synthesis-outlook}

The new constructions of this appendix fit together into a single type-$A$
picture.

\begin{center}
\begin{tabular}{c}
compact polynomial seed \\
$\Downarrow$ \\
weightwise completed RTT tower \\
$\Downarrow$ \\
dg-shifted Yangian via completed bar duality \\
$\Downarrow$ \\
KR telescopes and Baxter--Rees family \\
$\Downarrow$ \\
formal braided boundary germ \\
$\Downarrow$ \\
completed/coderived DK extension
\end{tabular}
\end{center}

At the compact end, the vector seed controls the Schur envelope and, in
the strict regime, the entire higher-point ordered packet. At the
completion end, the correct inverse-limit package is weightwise rather
than global. At the boundary end, the negative prefundamental object is
not an isolated module but the special fiber of an explicit affine family,
and not only the objects but also the braiding and associator data extend
across the boundary.

\begin{remark}[Relation with factorization quantum groups]
\label{rem:relation-factorization-quantum-groups}
Latyntsev shows that factorization quantum groups are controlled by
spectral $R$-matrices~\cite{Latyntsev23}. The continuation theorems of
this appendix can therefore be read as a boundary regularity principle for
the spectral side as well: once a compact equivalence matches the ordered
vector seed and the KR packet, the boundary tensor germ is no longer
additional mysterious structure but a continuation of the same spectral
package.
\end{remark}

\begin{remark}[Modular lift]
\label{rem:modular-lift-beyond-boundary}
The present appendix is genus zero in character, but it points directly
toward the modular lift. The monograph's universal class
$\Theta_{\mathcal A}$ already governs the modular deformation tower. Once
the genus-zero Baxter--Rees boundary germ and its tensor data are in
place, the next natural theorem is that the KZ monodromy of the compact
polynomial core deforms to a KZB-type modular monodromy controlled by the
same $\Theta_{\mathcal A}$ after passage to the completed boundary family.
That theorem is not proved here; what is proved is the algebraic and
representation-theoretic infrastructure it requires.
\end{remark}

\begin{remark}[Bibliographic note]
\label{rem:bibliographic-note-typeA-baxter-rees}
For integration into the repository, this appendix uses the existing keys
\texttt{DNP25}, \texttt{Latyntsev23}, \texttt{Zhang24},
\texttt{KTWWY14}, \texttt{BD04}, and \texttt{FG12}, together with the new
keys \texttt{HZ24} and \texttt{ZhangTheta24} supplied in the accompanying
bibliography patch.
\end{remark}

% ======================================================================
% WAVE-14 PLATONIC ANCHOR BLOCK (B3 upgrade, 2026-04-17)
% Locks AP159 four-Yangian-types disambiguation (this appendix is type-A
% specific, working in classical type-1 Y_h(sl_N) with type-2 dg-shifted
% completion), AP161 Notion-B working default, AP-RMATRIX level-prefix
% discipline for type-A R-matrices, AP170 equivalence of two chiral
% Yangian definitions OPEN, and Wave-14 cross-references to the unified
% Vol I Theorem A and Vol II climax. All labels new.
% ======================================================================

\section{Adversarial discipline block (Wave-14 anchor)}
\label{sec:typeA-baxter-rees-wave14-anchor}

\subsection{Type-A scope and the four-Yangian-types disambiguation}
\label{subsec:typeA-baxter-rees-ap159-yangian-types}

\begin{convention}[AP159: type-A appendix uses
$Y_\hbar(\mathrm{sl}_N)$ in classical (type-1) and dg-shifted
(type-2) presentations; \ClaimStatusProvedHere]
\label{conv:typeA-baxter-rees-ap159-yangian-types}
This appendix is genuinely type-A-specific (the seed rigidity
mechanism uses Schur--Weyl duality on $V = \C^N$ which fails in
non-type-A). In the four-Yangian-types disambiguation of Vol~I
AP159, the appendix uses:
\begin{enumerate}[label=\textup{(\arabic*)}]
\item \textbf{Classical $Y_\hbar(\mathrm{sl}_N)$} (type 1, $E_1$-top
on $\R$): the polynomial seed on $V\otimes V$, the rational
$R$-matrix from RTT, the asymptotic and prefundamental
modules; this is the level at which the seed propagation theorem
and weightwise stabilization are proved.
\item \textbf{dg-shifted $Y^{\mathrm{dg}}_\hbar(\mathrm{sl}_N)$}
(type 2, formal disk): the completed weightwise object built by
the Baxter--Rees family construction. Relating the spectral
factorization side
(type-4, Latyntsev~\cite{Latyntsev23}) to type 2 is the bridge of
\S\ref{chap:dg-shifted-factorization}.
\end{enumerate}
The chiral Yangian $Y(\mathrm{sl}_N)^{\mathrm{ch}}$ (type 3) appears
implicitly through the modular lift remark
(\S\ref{rem:modular-lift-beyond-boundary}): the type-A weightwise
completion of this appendix is the formal-disk shadow of the
type-3 chiral Yangian, and the lift to the curve $X$ requires
the modular comparison theorem (frontier).
\end{convention}

\begin{remark}[AP161 Notion-B for the type-A Baxter--Rees family;
\ClaimStatusProvedElsewhere]
\label{rem:typeA-baxter-rees-ap161-notion-b}
The Baxter--Rees family is constructed in Notion~B
($A_\infty$ in $\mathrm{End}^{\mathrm{ch}}_{Y_\hbar}$): the
weightwise stabilization of the bar complex
(prop:weightwise-stabilization) is a chain-level statement, the
strict regime
(thm:strict-rigidity, etc.) is the truncation $m_k=0$ for $k\ge 3$,
and the strictification beyond the strict regime
(higher Schur-functor towers) is the chain-level lift to nontrivial
$\{m_k\}_{k\ge 3}$. The orthogonal coideal mechanism of
\S\ref{chap:shifted-rtt-orthogonal-coideals} is the open-colour
Koszul-dual presentation in the same Notion-B language.
\end{remark}

\subsection{AP-RMATRIX guard for type-A spectral kernels}
\label{subsec:typeA-baxter-rees-rmatrix-level-prefix}

\begin{remark}[Level prefix on the type-A rational $R$-matrix;
AP-RMATRIX/HZ-1; \ClaimStatusProvedElsewhere]
\label{rem:typeA-baxter-rees-rmatrix-level-prefix}
The type-A rational $R$-matrix
$R^{\mathrm{rat}}(u) = 1 + (\hbar P)/u + O(\hbar^2)$ on
$V\otimes V$ ($V = \C^N$, $P$ the permutation, $u$ the spectral
parameter) carries the level prefix $\hbar$ implicitly. The
$\hbar=0$ specialisation gives $R^{\mathrm{rat}}(u) = 1$
(trivial $R$, abelian limit), matching the AP126 discipline:
substituting the level parameter to zero gives the trivial
$r$-matrix. The KZ-equivalent normalisation
$R^{\mathrm{KZ}}(u) = u/(u-\hbar)\cdot 1 + (\hbar)/(u-\hbar)\cdot P$
gives the same $\hbar=0$ trivial limit. The
classical limit $r(u) = -P/u$ at $\hbar\to 0$
(after extracting the $\hbar$ prefix) is the rational $r$-matrix
satisfying CYBE, with the level prefix $-1$ being the $\mathrm{sl}_N$
trace-form Casimir on $V\otimes V$.
\end{remark}

\begin{remark}[AP170 equivalence of two chiral Yangian definitions
OPEN even in type A]
\label{rem:typeA-baxter-rees-ap170-two-definitions}
Vol~I AP170: the two chiral Yangian definitions
(weaker $E_1$-factorization-plus-RTT vs stronger
modular-MC-tower) are equivalent at chain level only conditional
on the curved-Dunn $H^2 = 0$ at $g\ge 2$. Even in type A, where
the seed rigidity mechanism gives explicit control of the
$A_\infty$ tower, this equivalence is OPEN at chain level
beyond genus zero. The modular lift remark
(\S\ref{rem:modular-lift-beyond-boundary}) is the type-A
specialisation of this open question.
\end{remark}

\subsection{Wave-14 anchor: type-A bridge feeds the universal
holography master theorem}
\label{subsec:typeA-baxter-rees-wave14-anchor}

\begin{remark}[Type-A Baxter--Rees family as the
$\fg = \mathrm{sl}_N$ specialisation of the unified $Q_\fg^{k,f,\mu}$;
\ClaimStatusProvedElsewhere]
\label{rem:typeA-baxter-rees-wave14-uhm}
The Baxter--Rees family of this appendix supplies the type-A
specialisation of the unified chiral quantum group
$Q_\fg^{k,f,\mu}$
(Theorem~\ref{thm:unified-chiral-quantum-group}): for
$\fg = \mathrm{sl}_N$, $f = 0$ (unreduced), $\mu = 0$ (unshifted),
$k$ generic, the unified construction recovers the type-A Yangian
$Y_\hbar(\mathrm{sl}_N)$, and the Baxter--Rees compactification
of this appendix is the explicit weightwise model. The boundary
tensor germ
(\S\ref{rem:modular-lift-beyond-boundary}) is the input to
\cref{thm:universal-holography-master} (Vol~II
\cref{ch:programme-climax-platonic}) at the type-A line: the universal
holography functor $\Phi_{\mathrm{hol}}$ takes
$Y_\hbar(\mathrm{sl}_N)$-modules as Wilson lines and the bulk
$3$d HT theory $T_{Y_\hbar(\mathrm{sl}_N)}$ is the type-A holomorphic
Chern--Simons of Costello--Yamazaki (4d holomorphic--topological
mixed CS, restricted to $X\times\R$). The Vol~I unified Koszul
reflection (\cref{V1-thm:koszul-reflection}) controls the
type-A duality on the bar side; the universal conductor of
\cref{V1-chap:universal-conductor} controls the level-shift
$k\mapsto -k - 2N$ ($\mathrm{sl}_N$ critical level $-N$ via
$h^\vee = N$) on the conductor strata.
\end{remark>
